\paragraph{\trFukayaCategory。}\label{sub:the-fukaya-category}
或许有益的是回顾一下：\trAinfty{}结合性关系\eqref{eq:aoo-associativity} 之所以出现在\trFukayaCategory{}中，很大程度上要归功于\trDeligneMumfordStasheff{}\trCompactification{}$\bar{\scr{R}}^{d+1}$ 的递归结构。更确切地说，每一棵具有两个\trVertex{}的稳定 $d$ \trLeafedTree{}——按\cref{fig:two-vertex-tree}的记号即 $T(m,n)$——都标记着一个\trCodimension{}为一的\trBoundaryFace
\begin{equation}\label{eq:two-vertex-boundary-face}
\bar{\scr{R}}^{T(m,n)}\cong\bar{\scr{R}}^{d+2-m}\times\bar{\scr{R}}^{m+1}\subset\partial\bar{\scr{R}}^{d+1};
\end{equation}
这些\trBoundaryFace{}对应于\eqref{eq:aoo-associativity} 中 $1<m<d$ 的那些求和项。
正如从\eqref{eq:family-moduli-dimension} 所能看到的，\trAinfty{}结构的定义需要对诸空间 $\bar{\scr{R}}^{d+1}$ 的一个\trOrientation{}选取。一旦作出了该选取，\eqref{eq:two-vertex-boundary-face} 中的一次\trOrientation{}比较就对每一对 $(m,n)$ 给出一个符号。我们的目标是把事情安排成这样：这些符号，连同\eqref{eq:family-gluing-diagram} 中出现的那些 Koszul 符号、以及进入 $\mu^d$ 定义的一些其他辅助符号，最终凑出\eqref{eq:aoo-associativity} 的正确版本。我们从\eqref{eq:universal-disc-base} 出发，即把 $\scr{R}^{d+1}$ 定义为 $\mathrm{Conf}_{d+1}(\partial D)$ 对 $\mathrm{Aut}(D)$ 之商。取 $\mathrm{Conf}_{d+1}(\partial D)$ 中由使 $z_0$、$z_1$、$z_2$ 固定的构型 $(z_0,\ldots,z_d)$ 所组成的子空间。$(z_3,\ldots,z_d)\subset(\partial D)^{d-2}$ 的自然边界\trOrientation{}给出了该子空间的一个\trOrientation。另一方面，它是 $\mathrm{Aut}(D)$-作用的一个整体切片，从而可以经\trProjection{}与 $\scr{R}^{d+1}$ 等同。我们给 $\scr{R}^{d+1}$ 配上这一诱导的\trOrientation。一番容易但繁琐的计算（此处从略）表明：在\eqref{eq:two-vertex-boundary-face} 中，乘积\trOrientation{}$\bar{\scr{R}}^{d+2-m}\times\bar{\scr{R}}^{m+1}$ 与 $\partial\bar{\scr{R}}^{d+1}$ 的边界\trOrientation{}相差 $(-1)^{\ddagger}$，
\begin{equation}\label{eq:product-orientation-exponent}
\ddagger=m(d-n)+m+n.
\end{equation}
假设我们已经作好了搭建初步版本 $\scr{F}(M)^{\mathrm{pr}}$ 所需的全部选取，如\cref{sub:the-definition}中所解释的那样。定义真正的\trFukayaCategory{}$\scr{F}(M)$ 的头几步是显然的：\trObject{}是 $M$ 中的\trExactLagrangianBrane，而\trMorphism{}空间为
\begin{equation*}
\hom(L_0^{\#},L_1^{\#})=CF^*(L_0^{\#},L_1^{\#}),
\end{equation*}
一阶运算则是附加了一个符号的\trFloer{}\trBoundaryOperator，
\begin{equation*}
\mu^1(a)=(-1)^{|a|}\partial(a).
\end{equation*}
注意，这里用到的是与底层的\trLagrange\trSubmanifold{}对 $(L_0,L_1)$ 相关联的\trFloerDatum。原则上，也可以让\trFloerDatum（以及\trPerturbationDatum）同样依赖于诸\trBraneStructure；眼下这只会带来不必要的复杂化，但稍后在\trFukayaCategory{}的某些变体中我们将会采纳它，见 (14e)。其次，取\trBrane{}$(L_0^{\#},\ldots,L_d^{\#})$（对某个 $d\ge 2$），并取定 $y_0\in\scr{C}(L_0,L_d)$、$y_k\in\scr{C}(L_{k-1},L_k)$。利用\eqref{eq:family-moduli-dimension} 以及 $\scr{R}^{d+1}$ 的给定\trOrientation，可以看到诸\trModuliSpace{}$\scr{M}^{d+1}(y_0,\ldots,y_d)$ 的维数为 $i(y_0)-i(y_1)-\cdots-i(y_d)+2-d$，并且相对于 $o(y_0)\otimes o(y_1)^{\vee}\otimes\cdots\otimes o(y_d)^{\vee}$ 取\trOrientation。现在假设 $(r,u)$ 是某个零维\trModuliSpace{}中的一点。那么我们有 $i(y_0)=i(y_1)+\cdots+i(y_d)+2-d$，这正是运算 $\mu^d$ 所要求的\trGrading，此外我们还得到一个取定的\trIsomorphism
\begin{equation}\label{eq:preferred-isomorphism-mu}
c_{r,u}\colon o(y_d)\otimes\cdots\otimes o(y_1)\longrightarrow o(y_0).
\end{equation}
把 $\mu^d$ 的 $(y_0,\ldots,y_d)$-\trCoefficients{}定义为相应诸 $\KK$-\trNormalization{}$|c_{r,u}|_\KK$ 之和，一如\eqref{eq:boundary-operator-normalization} 那样，但要附加一个符号 $(-1)^{\dagger}$，
\begin{equation}\label{eq:dagger-sign-exponent}
\dagger=i(y_1)+2i(y_2)+\cdots+d\,i(y_d).
\end{equation}
我们要立刻指出：由于\eqref{eq:cohomological-composition} 中的约定，其底\trCohomology{}\trCategory{}确实与先前定义的 $H\scr{F}(M)$ 相同。
\begin{proposition}\label{thm:mu-form-a-infinity-structure}
诸 $\mu^d$ 构成一个\trAinfty{}结构。
\end{proposition}
\begin{proof}
正如前一小节中所做的那些较简单的例子一样，这基于一个\trGluing{}论证：从一对\trIsolated{}点 $(r_k,u_k)$（$k=1,2$）出发，如\eqref{eq:two-vertex-boundary} 那样。我们只考虑稳定情形 $1<m<d$，并有意识地忽略不稳定的情形（把该情形也纳入考虑所需的附加论证，紧随\eqref{eq:chain-map-boundary-formula} 与\eqref{eq:chain-homotopy-formula} 的证明之路，留作读者的练习）。取\trIsomorphism{}$c_{r_1,u_1}$ 与 $c_{r_2,u_2}$，把第二个的输出插入到第一个的第 $(n+1)$ 个输入中，使二者\trComposition。用\trOrientation{}空间的语言说，这对应于把 $o(\tilde{y})$ 与其对偶作\trPairing。为简洁计，我们把所得结果记作 $\tilde{c}$。由前文的一般性讨论，$\tilde{c}$ 诱导一维\trModuliSpace{}$\scr{M}^{d+1}(y_0,\ldots,y_d)$ 的一个\trOrientation。取某个小的参数 $\rho$，把诸 $(r_k,u_k)$ \trGluing{}成该空间中的一点 $(r,u)$。按\eqref{eq:polygon-compactification} 的\trGromov{}\trCompactification{}的语言，$(r,u)$ 紧挨着边界点 $(r_1,u_1,r_2,u_2)$，因而存在一个指向该边界的\trCanonical{}局部\trOrientation。这两个\trOrientation{}相差一个符号 $(-1)^{\text{\S}}$。一旦确定了该符号，由构造即可推知：把诸\trNormalization{}$(-1)^{\text{\S}}\tilde{c}|_\KK$ 对一切对 $(r_k,u_k)$ 加总，其和为零。要由此过渡到 $(r_1,u_1,r_2,u_2)$ 对 $\mu^{d+1-m}(a_d,\ldots,\mu^m(a_{n+m},\ldots,a_{n+1}),a_n,\ldots,a_1)$ 的 $(y_0,\ldots,y_d)$-\trCoefficients{}所作的贡献，还须补上来自\eqref{eq:dagger-sign-exponent} 的相应符号，
\begin{equation*}
\begin{aligned}
\dagger_1&=i(y_1)+2i(y_2)+\cdots+(n+1)i(\tilde{y})\\
&\quad+(n+2)i(y_{n+m+1})+\cdots+(d+1-m)i(y_d)\\
&=i(y_1)+2i(y_2)+\cdots+(n+1)\bigl(i(y_{n+1})+\cdots+i(y_{n+m})+2-m\bigr)\\
&\quad+(n+2)i(y_{n+m+1})+\cdots+(d+1-m)i(y_d),\\
\dagger_2&=i(y_{n+1})+2i(y_{n+2})+\cdots+m\,i(y_{n+m}).
\end{aligned}
\end{equation*}
余下的断言是
\begin{equation}\label{eq:sign-identity-clubsuit}
\text{\S}+\dagger_1+\dagger_2=\clubsuit_n+\bigl(2i(y_1)+\cdots+(d+1)i(y_d)\bigr),
\end{equation}
其中右端第一项给出\eqref{eq:aoo-associativity} 的标准符号，而第二项与 $(m,n)$ 无关。这与上面的论证合在一起，便确立了各结合性方程。
为完成必要的计算，考虑图式\eqref{eq:family-gluing-diagram}，取 $\scr{R}_1=\scr{R}^{d+2-m}$、$\scr{R}_2=\scr{R}^{m+1}$，但现在经由与 $T(m,n)$ 相关联的映射\eqref{eq:gluing-map}，把 $\scr{R}=(-1;0)\times\scr{R}_1\times\scr{R}_2$ 换成 $\scr{R}^{d+1}$。相应地，我们令 $\scr{M}_{\scr{S}_1}=\scr{M}^{d+2-m}$、$\scr{M}_{\scr{S}_2}=\scr{M}^{m+1}$，并把 $\scr{M}_{\scr{S}}$ 换成 $\scr{M}^{d+1}$。在图式的每个角上都坐着一个一维实\trVectorSpace，我们按如下方式为它们配置\trOrientation。
\begin{itemize}
\item 在右上角，取 $\tilde{c}$ 连同 $\scr{R}^{d+1}$ 的\trOrientation，然后用\trHorizontal{}箭头把它\trTransfer{}到 $\lambda^{\mathrm{top}}T\scr{M}^{d+1}$ 的一个元素上。用我们先前的讨论来说，这恰恰就是由 $\tilde{c}$ 诱导的该\trModuliSpace{}的\trOrientation。
\item 在左下角，该\trVectorSpace{}被\trCanonical{}地等同于 $\RR$，我们用\trHorizontal{}箭头把它\trTransfer{}到右下角的一个\trTrivialization{}上。按诸 $c_{r_k,u_k}$ 的定义，这与取 $\scr{R}^{d+2-m}$ 和 $\scr{R}^{m+1}$ 的\trOrientation{}并把它们与 $c_{r_1,u_1}\otimes c_{r_2,u_2}$ 组合起来是同一回事。
\end{itemize}
其次，考虑这些\trOrientation{}是否与图式中的\trVertical{}箭头\trCompatible。
\begin{itemize}
\item 若取左下角诸\trModuliSpace{}的平凡\trOrientation，并经 $\uparrow$——即\trGluingMap{}的\trLinearization——把它映过去，便得到指向\trGromov{}边界的 $\scr{M}^{d+1}$ 的一个\trOrientation，按定义它与上面引入的\trOrientation{}相差 $(-1)^{\text{\S}}$。
\item 在右列中，两个\trOrientation{}之间仅有的差别来自诸相关\trFamily{}的底空间之间的\trGluingMap{} \eqref{eq:gluing-map}，因而等于 $(-1)^{\ddagger}$。
\end{itemize}
最后，整幅图式在相差一个已知的 Koszul 符号 $(-1)^{\circ}$ 的意义下交换，在此情形 $\circ=m(d-m-1)+m\sum_{k>n+m}i(y_k)$。于是我们知道 $\text{\S}=\circ+\ddagger\pmod{2}$。一番直截了当的计算便给出\eqref{eq:sign-identity-clubsuit}。
\end{proof}
\begin{remark}\label{thm:sign-convention-remark}
我们的构造中有两个人为的要素：其一是 $\bar{\scr{R}}^{d+1}$ 的诸\trOrientation（或者更确切地说，是要求诸\trBoundaryStratum{}给出\trOrientation{}差\eqref{eq:product-orientation-exponent} 的那些要求——它们在相差一个无关宏旨的整体变号的意义下确定那些\trOrientation），其二是\eqref{eq:dagger-sign-exponent}。通过调整这两项，可以得到当前通行的任何一种\trAinfty{}结构符号约定；几何本身并不挑出任何优先的约定。\cref{thm:mu-form-a-infinity-structure}的另一个证明见 [60, \S22]。
\end{remark}
